\section{程序源代码}\label{app:code}

以下为本文全部计算所用的源代码，语言为 Python，依赖仅限 NumPy、SciPy 与 openpyxl。运行顺序为参数与数据核对、四问求解、校验与对比、结果导出，由 \texttt{main.py} 统一调度。文中所有数值均由这些代码一次运行产出，未经手工修改。

\subsection{参数与题面数据}

\lstinputlisting[caption={params.py：题面原值的唯一存放处}]{../code/params.py}

\lstinputlisting[caption={gen\_params.py：由附件生成时变边界与半径序列}]{../code/gen_params.py}

\lstinputlisting[caption={data\_check.py：附件数据的完整性与单调性核对}]{../code/data_check.py}

\subsection{公共求解工具}

\lstinputlisting[caption={utils.py：物性公式、界面积分平均、控制体体积与三对角求解}]{../code/utils.py}

\subsection{四问的求解}

\lstinputlisting[caption={problem1.py：预热阶段常物性耦合场}]{../code/problem1.py}

\lstinputlisting[caption={problem2.py：全过程变物性耦合场}]{../code/problem2.py}

\lstinputlisting[caption={problem3.py：逐点判据下的首次达标停时}]{../code/problem3.py}

\lstinputlisting[caption={problem4.py：收缩域上的移动边界模型（含实验室坐标对照支）}]{../code/problem4.py}

\subsection{校验、对比与灵敏度}

\lstinputlisting[caption={verification.py：网格与时间步收敛、守恒残差、解析解对照、退化算例}]{../code/verification.py}

\lstinputlisting[caption={model\_compare.py：集中参数与常物性两个候选模型的实算}]{../code/model_compare.py}

\lstinputlisting[caption={sensitivity.py：基准与十三支参数扰动、建模选择对照算例}]{../code/sensitivity.py}

\lstinputlisting[caption={constraint\_audit.py：结果文件规格与题面要求的逐项核对}]{../code/constraint_audit.py}

\subsection{结果导出与调度}

\lstinputlisting[caption={export\_excel.py：四份结果文件的生成}]{../code/export_excel.py}

\lstinputlisting[caption={main.py：全流程调度}]{../code/main.py}
